\documentclass[
    language=japanese,
    doctype=article,
    institution=none,
]{../../omnilatex}

\RequirePackage{config/document-settings}

\begin{document}

\maketitle

\section*{概要}
この文書は、OmniLaTeX の日本語組版機能を示します。日本語技術文書の作成を通じて、字体設定、行間処理、ルビ注釈などの機能を紹介します。

\tableofcontents

\section{はじめに}
OmniLaTeX は LuaLaTeX と luatexja エンジンに基づき、日本語文書のネイティブな組版を提供します。

日本語の組版は中国語や韓国語とは異なる多くの特徴を持っています。例えば、行頭行末の禁則処理（句読点の配置規則）、ルビ（ふりがな）の自動配置、括弧類の自動置換などです。これらの処理を正しく行うことが、読みやすく美しい日本語文書を作成する上で不可欠です。

近年、日本のソフトウェア業界では国際化対応の重要性が増しています。日本語ドキュメントの品質向上は、製品の信頼性と使いやすさに直結します。

\section{日本語組版の機能}

\subsection{基本組版}
吾輩は猫である。名前はまだ無い。どこで生れたかとんと見当がつかぬ。
何でも薄暗いじめじめした所でニャーニャー泣いていた事だけは記憶している。

日本語の段落は原則として一字下げ（字下げ）で始まります。行間は通常、字号の 1.5 倍から 2 倍に設定されます。句読点は全角で表記され、OmniLaTeX は自動的に行末の句点と行頭の開き括弧の処理を行います。

\subsection{和英混植}
日本語と英語が混在する文書では、両者の間に適切な空白を挿入する必要があります。例えば、「アルゴリズムの計算量は $O(n \log n)$ である」という文では、「$O$」の前後に自然な間隔が必要です。

OmniLaTeX は luatexja の機能を利用して、和文字と欧文字の間に自動的に四分アキ（文字の四分の幅の空白）を挿入します。この処理はフォントのグリフ幅に基づいて動的に計算されるため、手動の調整は不要です。

\subsection{ルビ（ふりがな）}
OmniLaTeX はルビ注釈をサポートします：

%furigana{漢}{かん}%furigana{字}{じ}の%furigana{読}{よ}み%furigana{方}{かた}

ルビは主に漢字の読み方を示すために使用されますが、専門用語に略称を付記する用途にも使われます。例えば、アプリケーション・プログラミング・インターフェースを「%furigana{API}{エーピーアイ}」と表記することができます。ルビのフォントサイズは親文字の半分程度に設定されるのが一般的です。

\subsection{縦書きモード}
伝統的な日本語文書では縦書きが使われます：

% \begin{vertical} — requires luatexja-patch; disabled for TL2025 compat
% 吾輩は猫である。名前はまだ無い。どこで生れたか見当がつかぬ。
% \end{vertical}

縦書きモードは小説や詩歌などの文芸作品で広く使用されます。技術文書では横書きが主流ですが、表紙や章扉のデザインに縦書きを取り入れることで、和風の美しさを演出することができます。

\section{日本語技術文書の実践}

\subsection{仕様書の書き方}
ソフトウェア仕様書は、システムの機能要件と非機能要件を明確に記述する文書です。日本の業界では JIS X 0131 に準拠した仕様書の書き方が推奨されています。OmniLaTeX を使用することで、見出し番号の自動生成、相互参照、用語集の管理などの機能を活用でき、仕様書の作成効率が大幅に向上します。

\subsection{論文の構成}
日本語の学術論文は、概要、序論、本論、結論の四部分構成が一般的です。数式を含む理工系論文では、$E = mc^2$ のような行内数式や、独立した行間数式の正確な組版が不可欠です。OmniLaTeX は \LaTeX{} の強力な数式組版機能を日本語文脈でそのまま利用できます。

\subsection{数式の例}
日本語の数学・理工系文書では、複雑な数式の組版が頻繁に必要になります。例えば、オイラーの等式は $e^{i\pi} + 1 = 0$ と表記され、ガウス積分は以下のように組版できます：
\[
\int_{-\infty}^{+\infty} e^{-x^2} \, dx = \sqrt{\pi}
\]

\subsection{開発ドキュメントの国際化}
日本の開発チームがグローバルに展開する際、日本語と英語の両方でドキュメントを整備する必要があります。OmniLaTeX の多言語対応機能により、同じテンプレートを用いて言語を切り替えることが可能です。これにより、ドキュメントの構成やスタイルの一貫性を保ちながら、複数言語の版を効率的に作成できます。

\section{まとめ}
OmniLaTeX は日本語の学術文書に完全な組版ソリューションを提供します。LuaLaTeX と luatexja の統合により、ルビ注釈、禁則処理、和英混植の間隔調整など、日本語特有の組版要件を自動的に処理します。技術文書から文芸作品まで、幅広いジャンルの日本語ドキュメントに対応可能です。

\end{document}
